O trabalho foi conduzido sob a liderança de \textbf{Suhas Parandekar}, Economista Sênior. 
A equipe técnica foi composta por \textbf{Paulo Nascimento}, responsável pela dimensão de demanda de cursos técnicos; \textbf{Gabriel Grabowski}, responsável pela frente de oferta da Educação Profissional e Tecnológica (EPT); \textbf{Fábio Maciel}, responsável pelo desenvolvimento e análise de dados; e 
\textbf{Marina Bomura}, responsável pela coordenação e gestão do projeto.

Este relatório se beneficiou de discussões com outros membros da equipe do Banco Mundial para educação no Brasil, entre eles \textbf{Paulo Facundo Cuevas}, Economista-Líder; \textbf{Leandro Oliveira Costa}, Economista Sênior; \textbf{Ildo José Lautharte Júnior}, Economista Sênior; \textbf{Louisee Rodrigues da Cruz Boari}, Consultora Especialista em Educação; e \textbf{Isabella Miranda Meyer}, Associada Profissional Júnior.  O apoio administrativo foi prestado por \textbf{Débora Calheiros Fortes}.

A equipe trabalhou sob a orientação de \textbf{Andreas Blom}, Gerente de Prática de Educação do Banco Mundial para a Região da América Latina e Caribe.  
A \textbf{Diretora do Banco Mundial para o Brasil} é \textbf{Cécile Fruman}, e o \textbf{Diretor Global de Educação} é \textbf{Jaime Saavedra}.

A equipe é extremamente grata a \textbf{José Henrique Paim Fernandez}, Diretor da \textbf{Fundação Getulio Vargas – Diretoria de Desenvolvimento da Gestão Pública e Políticas Educacionais}, e à equipe da 
\textbf{FGV-DGPE} pela valiosa colaboração analítica e técnica.
Agradecemos, em especial, a \textbf{Aléssio Trindade} pelo apoio estratégico e pelas contribuições substanciais ao desenvolvimento do trabalho;  
a \textbf{Tullia Maria Erseni} pela coordenação e parceria constante;  
a \textbf{Carlos Ávila} e \textbf{Rogério Vidal} pelo suporte analítico na área de dados financeiros;  
e a \textbf{Garabed Kenchian} pelo apoio na análise de dados relacionados à Educação Profissional e Tecnológica (EPT).

A equipe expressa seu agradecimento ao \textbf{Ministério da Educação}, em especial à \textbf{Secretaria de Educação Profissional e Tecnológica (SETEC)} pelas conversas inicias sobre o PROPAG e ao \textbf{Instituto Nacional de Estudos e Pesquisas Educacionais Anísio Teixeira (MEC/INEP)}, bem como ao \textbf{Instituto Brasileiro de Geografia e Estatística (IBGE)} e ao \textbf{Instituto de Pesquisa Econômica Aplicada (IPEA)}, pelo fornecimento de informações estatísticas e socioeconômicas fundamentais para a análise e elaboração deste estudo.


\vspace{2em}

\textbf{Contactos:} sparandekar [at]  worldbank [dot]  org; 



\textbf{\textcolor{orange}{Plataforma Interativa na Internet:}}

\href{https://datanalytics.worldbank.org/content/6604f317-d487-436a-b6b6-792369b341f0}{https://datanalytics.worldbank.org/Plataforma Interativo PROPAG} 
